%*******************************************************
% Zusammenfassung
%*******************************************************
%\renewcommand{\abstractname}{Abstract}
\begin{otherlanguage}{ngerman}
\pdfbookmark[1]{Zusammenfassung}{Zusammenfassung}
\begingroup
\let\clearpage\relax
\let\cleardoublepage\relax
\let\cleardoublepage\relax

\chapter*{Zusammenfassung}
In der strukturellen Beweistheorie kombiniert der Kalkül des natürlichen
Schließens ($\ND$) Einführungs- und Eliminationsregeln, während Gentzens
$\LK$ durch duale Einführungsregeln definiert ist.
Durch die Curry-Howard-Korrespondenz erfassen der $\lambda$-Kalkül
und der $\lambda\mu\tilde{\mu}$-Kalkül den rechnerischen Kern der Normalisierung
in diesen jeweiligen Systemen. Es existieren jedoch auch alternative logische Systeme
wie Parigots \enquote{Free Deduction} ($\FD$), welche
ausschließlich auf Eliminationsregeln basieren.
Zwar lassen sich das Termzuweisungssystem und die Semantik eines solchen Kalküls
anhand seiner Übersetzung in entweder $\ND$ oder $\LK$ verstehen, doch fehlt
uns eine eigenständige operationalle Interpretation für
ein solches System.

In dieser Arbeit präsentieren wir mit $\piFD$ eine moderne Rekonstruktion von Free Deduction
als Prozesskalkül auf der Basis linearer Logik---ähnlich zu Negris
\enquote{Uniform Calculus for Classical Linear Logic} (UCL).
Indem wir logische Axiome als Session-Kanäle und Eliminationsregeln
als die Erfüllung eines Futures bzw. Promises interpretieren,
erhalten wir einen Kalkül, in dem die Schnittregel (\textsc{cut}) als
Fork-Primitive dient, die lediglich die Interaktion zwischen parallelen Prozessen,
nicht jedoch die Kommunikation zwischen Endpunkten eines Kommunikationskanals ermöglicht.
Dieser Ansatz liefert eine alternative Grundlage für den Entwurf von
Session-Typ-basierten Sprachen, die von den traditionellen
Curry-Howard-Interpretationen von Girards linearer Logik abweicht.

Wir formalisieren die Metatheorie von $\piFD$ in Rocq und liefern maschinengeprüfte
Beweise für die Typsicherheit und die Church-Rosser-Eigenschaft.
Diese Ergebnisse schaffen eine fundierte Grundlage für die Analyse
kommunizierender, nebenläufiger Prozesse in eliminationsbasierten Systemen.
Letztendlich vertieft diese Arbeit die Curry-Howard-Korrespondenz für
parallele und kommunikationsbasierte Kalküle und bietet neue Einblicke
in die Beziehung zwischen logischer Elimination und Session-basierter
Nebenläufigkeit.

\end{otherlanguage}

\endgroup

\vfill